\documentclass[10pt]{article}
\usepackage[T1]{fontenc}
\usepackage[utf8]{inputenc}
\usepackage[french]{babel}
\usepackage{amsmath,amssymb,mathrsfs,mathpazo}
\usepackage{geometry}
\geometry{a4paper,margin=2cm}
\usepackage{graphicx,array,tabularx,tabularray,colortbl,multirow,multicol,booktabs,enumitem,pifont,ulem,fancybox,diagbox,cancel}
\usepackage{pstricks,pstricks-add,pst-plot,pst-tree,pst-node,pst-text,pst-func,pst-math,pst-bspline,pst-3dplot,pst-eucl,pst-solides3d}
\let\pscalculate\undefined
\usepackage{pst-bezier}
\usepackage{tikz,pgfplots,tkz-tab}
\usetikzlibrary{arrows,arrows.meta,patterns,calc,positioning,trees,decorations.markings,shapes.geometric,shapes.misc,shapes.arrows,intersections,decorations.pathmorphing,decorations.pathreplacing,backgrounds,shadows,babel}
\pgfplotsset{compat=1.18}
\usepgfplotslibrary{fillbetween}
\usepackage[np]{numprint}
\usepackage{listings,makecell,pas-tableur,twoopt}
\let\barre\undefined
\usepackage[pstricks]{tablvar}
\definecolor{darkgreen}{RGB}{0,100,0}
\providecommand{\red}{\color{red}}
\providecommand{\blue}{\color{blue}}
\providecommand{\green}{\color{green}}
\providecommand{\black}{\color{black}}
\providecommand{\white}{\color{white}}
\providecommand{\gray}{\color{gray}}
\providecommand{\cyan}{\color{cyan}}
\providecommand{\magenta}{\color{magenta}}
\providecommand{\yellow}{\color{yellow}}
\providecommand{\violet}{\color{violet}}
\providecommand{\orange}{\color{orange}}
\providecommand{\euro}{\texteuro}
\providecommand{\EUR}{\texteuro}
\providecommand{\ieme}{\textsuperscript{e}}
\providecommand{\degres}{\ensuremath{^\circ}}
\providecommand{\degree}{\ensuremath{^\circ}}
\providecommand{\vv}[1]{\overrightarrow{#1}}
\providecommand{\Coord}[1]{#1}
\providecommand{\decosix}{}
\providecommand{\starredbullet}{\ensuremath{\bullet}}
\pagestyle{empty}
\setlength{\parindent}{0pt}
\usepackage{adjustbox}\DeclareUnicodeCharacter{2236}{:}\DeclareUnicodeCharacter{2190}{\ensuremath{\leftarrow}}
\providecommand{\R}{\mathbb{R}}
\providecommand{\N}{\mathbb{N}}
\providecommand{\Z}{\mathbb{Z}}
\providecommand{\Q}{\mathbb{Q}}
\providecommand{\D}{\mathbb{D}}
\providecommand{\C}{\mathbb{C}}
\providecommand{\e}{{\mathrm{e}}}
\providecommand{\vect}[1]{\overrightarrow{#1}}
\providecommand{\vectt}[1]{\overrightarrow{\mathrm{#1}}}
\providecommand{\vvt}[1]{\overrightarrow{\mathrm{#1}}}
\providecommand{\np}[2][]{{#2~#1}}
\providecommand{\vv}[1]{{\overrightarrow{#1}}}
\renewcommand{\d}{{\mathrm{d}}}
\providecommand{\degres}{{^\circ}}
\providecommand{\degre}{{^\circ}}
\providecommand{\blue}{\color{blue}}
\providecommand{\red}{\color{red}}
\providecommand{\enskip}{\quad}
\providecommand{\llbracket}{\lbrack\!\lbrack}
\providecommand{\rrbracket}{\rbrack\!\rbrack}
\providecommand{\euro}{{\text{€}}}
\providecommand{\no}{{\text{n°}}}
\providecommand{\interval}[2]{{[#1;#2]}}
\providecommand{\intervFO}[2]{{[#1;#2[}}
\providecommand{\dotfill}{\ldots}

\providecommand{\ds}{\displaystyle}
\providecommand{\textding}[1]{#1}
\providecommand{\Oij}{\(\mathrm{O};\vec{\imath},\vec{\jmath}\)}
\providecommand{\Oijk}{\(\mathrm{O};\vec{\imath},\vec{\jmath},\vec{k}\)}
\providecommand{\Ouv}{\(\mathrm{O};\vec{u},\vec{v}\)}
\providecommand{\intervFF}[2]{\ensuremath{[#1;#2]}}
\renewcommand{\euro}{\ifmmode\text{\texteuro}\else\texteuro\fi}
\DeclareUnicodeCharacter{00B0}{\ifmmode^{\circ}\else\textdegree\fi}
\DeclareUnicodeCharacter{20AC}{\ifmmode\text{\texteuro}\else\texteuro\fi}
% Énoncés : Éducation nationale. Transcription et corrigés : APMEP, Denis Vergès et contributeurs. Découpage et classement : CoopMaths, Sébastien Lozano, Éric Elter et contributeurs. Adaptation de la présentation : Maths Explorer.
% CC BY-SA : https://coopmaths.fr/www/about/
% https://coopmaths.fr/alea/static/bac/2024/tex/bac_2024_01_sujet0_6_cor.tex
\begin{document}
\noindent\textit{APMEP / CoopMaths. CC BY-SA. Présentation adaptée par Maths Explorer.}\par\medskip
On considère la suite $\left(u_{n}\right)$ définie par :
\[\left\{\begin{array}{l !{=} l}
u_{n + 1} & \dfrac{u_{n}}{1 + 2 u_{n}} \text { pour tout entier naturel } n 
\\ u_{0} & 1
\end{array}\right.\]

\begin{list}{\textbullet}{}
\item \textbf{Affirmation 1 :} \og{} $u_{4} = \dfrac{1}{9}$. \fg{}

$u_0=1$; 
$u_1=\dfrac{u_0}{1+2u_0} = \dfrac{1}{1+2\times 1}=\dfrac{1}{3}$;
$u_2=\dfrac{u_1}{1+2u_1} = \dfrac{\frac{1}{3}}{1+2\times \frac{1}{3}}
= \dfrac{\frac{1}{3}}{\frac{5}{3}} = \dfrac{1}{5}$;\\
$u_3=\dfrac{u_2}{1+2u_2} = \dfrac{\frac{1}{5}}{1+2\times \frac{1}{5}}
= \dfrac{\frac{1}{5}}{\frac{7}{5}} = \dfrac{1}{7}$;
$u_4=\dfrac{u_3}{1+2u_3} = \dfrac{\frac{1}{7}}{1+2\times \frac{1}{7}}
= \dfrac{\frac{1}{7}}{\frac{9}{7}} = \dfrac{1}{9}$

\hfill\textbf{Affirmation 1 vraie}

\item \textbf{Affirmation 2 :} \og{}Pour tout entier naturel $n$, $u_{n} = \dfrac{1}{2 n + 1}$.\fg{}

D'après la question précédente, la propriété \og $u_{n} = \dfrac{1}{2 n + 1}$ \fg{} est vérifiée pour $n$ entre 0 et 4. On va démontrer par récurrence qu'elle est vraie pour tout $n$.

\begin{list}{\textbullet}{}
\item \textbf{Initialisation}

Pour $n=0$, on a: $u_0=1$ et $\dfrac{1}{2\times 0 +1}=1$; donc la propriété est vérifiée.

\item On suppose la propriété vraie à un rang $n\geqslant 0$, c'est-à-dire que $u_n=\dfrac{1}{2n+1}$. C'est l'hypothèse de récurrence.

$u_{n+1}= \dfrac{u_n}{1+2u_n}
= \dfrac{\frac{1}{2n+1}}{1+2\frac{1}{2n+1}}
= \dfrac{\frac{1}{2n+1}}{\frac{2n+1+2}{2n+1}}
= \dfrac{1}{2n+1}\times \dfrac{2n+1}{2n+3}
=\dfrac{1}{2(n+1)+1}$

Donc la propriété est vraie au rang $n+1$.

\item \textbf{Conclusion}

La propriété est vraie au rang $0$, et elle est héréditaire pour tout $n\geqslant 0$, donc, d'après le principe de récurrence, elle est vraie pour tout $n$.
\end{list}

\hfill\textbf{Affirmation 2 vraie}

\item \textbf{Affirmation 3 :} \og{}La suite numérique $\left(u_{n}\right)$ est minorée par $10^{-10}$.\fg{}

$\ds\lim_{n\to +\infty} 2n+1=+\infty$ donc $\ds\lim_{n\to +\infty} \dfrac{1}{2n+1}=0$ et donc $\ds\lim_{n\to +\infty} u_n=0$.

Si la suite $(u_n)$ est minorée par $10^{-10}$ qui est un nombre strictement positif, elle ne peut pas avoir pour limite le nombre 0.

\hfill\textbf{Affirmation 3 fausse}

\end{list}

\end{document}
